\documentclass[a4paper,12pt]{article}

% Тот же набор средств оформления, что в математическом черновике проекта.
\usepackage{iftex}
\ifPDFTeX
  \errmessage{Для этого документа используйте XeLaTeX или LuaLaTeX}
\else
  \usepackage{fontspec}
  \usepackage{polyglossia}
  \setmainlanguage{russian}
  \IfFontExistsTF{DejaVu Serif}{\setmainfont{DejaVu Serif}}{\setmainfont{Times New Roman}}
  \IfFontExistsTF{DejaVu Sans Mono}{\newfontfamily\cyrillicfonttt{DejaVu Sans Mono}}{\newfontfamily\cyrillicfonttt{Menlo}}
\fi
\usepackage{amsmath,amssymb,mathtools}
\usepackage{booktabs,geometry,indentfirst,microtype,ragged2e,url}
\geometry{left=2.5cm,right=2.5cm,top=2.1cm,bottom=2.1cm}
\setlength{\parindent}{0.7cm}
\setlength{\parskip}{0pt}
\emergencystretch=4em
\tolerance=3000

\title{Реализованная модель и исходные данные\\
  \large Демонстрация первых попыток подбора пассивных орбит наблюдаемой сетки КА}
\author{И.~В.~Шимановская}
\date{Исследовательская версия, сентябрь 2026 г.}

\begin{document}
\justifying
\maketitle

\section{Назначение и граница результата}

Программа показывает один выбранный пролёт группы из $S=N^2$ космических
аппаратов (КА), $N=2,\ldots,10$. Цель исследования --- найти такие
\emph{начальные} состояния независимых пассивных орбит, при которых
вся группа некоторое время видна наземному наблюдателю как квадратная
сетка. Стартовый случай имеет $N=10$ ($S=100$) и целевую внешнюю
угловую сторону $20^\circ$. Вычисленный интервал одного пролёта есть
первый результат ограниченного поиска, а не проект готовой миссии,
не доказательство повторяемости и не прогноз предотвращения столкновений.
После общей эпохи аппаратам не назначается тяга и их координаты
не привязываются к идеальной решётке.

\section{Исходные данные и системы координат}

\subsection{Группа, орбита, наблюдатель и время}

Каждый КА рассматривается как пример платформы формата 16U.
Приняты $m=22$~кг, размеры корпуса
$0{,}454\times0{,}2463\times0{,}2463$~м,
эффективная площадь для сопротивления $A=0{,}1118202$~м$^2$,
$C_D=2{,}2$ и ориентация широкой гранью перпендикулярно мгновенной
скорости относительно атмосферы. Масса и габариты опубликованы для
\emph{другого} аппарата Corvus-HD 16U; для исследуемой группы они,
как и $A$, $C_D$ и ориентация, являются допущениями, а не
подтверждёнными параметрами изготовленного КА~\cite{corvus}.
Изготовленные размеры и характеристики двигателя этой группы неизвестны.

Общая эпоха элементов $t=0$ --- \texttt{2026-09-28T00:00:00Z}.
Опорный КА~0 на эпохе задан высотой 700~км и наклонением
$97{,}8^\circ$; его движение после эпохи вычисляется тем же
интегратором, что и для остальных КА. Наблюдатель расположен на
сферической Земле при геоцентрических широте
$\varphi=25{,}3026^\circ$ и долготе $\ell=51{,}5132^\circ$.
Приняты $R_E=6371$~км, $\mu=398600{,}4$~км$^3$/с$^2$,
$J_2=0{,}00108262668$ и
$\omega_E=7{,}292115\cdot10^{-5}$~рад/с.

Абсолютные радиусы и скорости задаются в геоцентрической
приближённой инерциальной системе ECI: ось $Z$ направлена вдоль оси
вращения Земли, ось $X$ приближённо связана со средним равноденствием
даты. Угол Гринвича на эпохе равен
$\Theta_0=6{,}785648977^\circ$ по формуле GMST при UT1$\approx$UTC.
Полное преобразование ITRF/ICRF, прецессия, нутация, полярное
движение и фактическая поправка UT1--UTC не используются.
Следовательно, это не точная эфемеридная ориентация на календарную дату.
Положение неподвижного наземного наблюдателя в той же ECI равно
\begin{equation}
 \mathbf R(t)=R_E
 \begin{pmatrix}
 \cos\varphi\cos\Lambda(t)\\
 \cos\varphi\sin\Lambda(t)\\
 \sin\varphi
 \end{pmatrix},\qquad
 \Lambda(t)=\ell+\Theta_0+\omega_E t.
 \label{eq:observer}
\end{equation}

Индексная матрица $N\times N$ постоянно связывает аппарат со строкой
и столбцом; в соседних строках порядок индексов чередуется.
При $N=2$ она равна $\bigl(\begin{smallmatrix}0&1\\3&2\end{smallmatrix}\bigr)$,
так что крайние вершины идут в порядке $0\!\to\!1\!\to\!2\!\to\!3$.
КА~0 задаёт направление локальных осей, но начало локальной системы
находится в \emph{геометрическом центре всех} $S$ КА:
\begin{equation}
 \mathbf r_c=\frac1S\sum_{j=0}^{S-1}\mathbf r_j,\qquad
 \mathbf v_c=\frac1S\sum_{j=0}^{S-1}\mathbf v_j,\qquad
 \mathbf e_x=\frac{\mathbf r_0}{|\mathbf r_0|},\quad
 \mathbf e_z=\frac{\mathbf r_0\times\mathbf v_0}
                       {|\mathbf r_0\times\mathbf v_0|},\quad
 \mathbf e_y=\mathbf e_z\times\mathbf e_x.
 \label{eq:center}
\end{equation}
Если $B=[\mathbf e_x\ \mathbf e_y\ \mathbf e_z]$, то
\begin{equation}
 \boldsymbol\rho_j=B^T(\mathbf r_j-\mathbf r_c),\qquad
 \dot{\boldsymbol\rho}_j=B^T(\mathbf v_j-\mathbf v_c)
                +\dot B^T(\mathbf r_j-\mathbf r_c).
 \label{eq:local}
\end{equation}
При вычислении $\dot B$ учитывается полное ускорение КА~0 из
уравнения~\eqref{eq:motion}. Центр --- расчётная точка, не
дополнительный аппарат и не самостоятельная кеплерова орбита.
Проекция наблюдателя строится из абсолютных $\mathbf r_j$ и
\eqref{eq:observer}, поэтому перенос локального начала её не меняет.

\section{Фактически используемая динамика и элементы орбит}

Каждый КА независимо интегрируется методом Рунге--Кутты четвёртого
порядка с максимальным шагом 2~с. Внутренние единицы --- км, с и
км/с. Реально включены \emph{только} центральное притяжение Земли,
вторая зональная гармоника $J_2$ и атмосферное сопротивление:
\begin{align}
 \ddot{\mathbf r}_j&=\mathbf g_{J_2}(\mathbf r_j)
                     +\mathbf a_D(\mathbf r_j,\mathbf v_j),
                     \label{eq:motion}\\
 \mathbf g_{J_2}(\mathbf r)&=-\frac{\mu\mathbf r}{r^3}
 +\frac{3J_2\mu R_E^2}{2r^5}
 \begin{pmatrix}
 x(5z^2/r^2-1)\\y(5z^2/r^2-1)\\z(5z^2/r^2-3)
 \end{pmatrix},\qquad r=|\mathbf r|,\label{eq:j2}\\
 \mathbf v_{\rm rel}&=\mathbf v-\boldsymbol\omega_E\times\mathbf r,
 \qquad \boldsymbol\omega_E=(0,0,\omega_E)^T,\nonumber\\
 \mathbf a_D&=-10^3\frac{\rho(h)C_DA}{2m}
 |\mathbf v_{\rm rel}|\mathbf v_{\rm rel},\qquad
 h=r-R_E,\label{eq:drag}\\
 \rho(h)&=4{,}6\cdot10^{-13}
 \exp\!\left(-\frac{h-700}{97}\right)\;\mathrm{kg/m^3}.
 \label{eq:density}
\end{align}
В~\eqref{eq:drag} $\rho,A,m$ даны в единицах СИ, скорость --- в
км/с; множитель $10^3$ возвращает ускорение в км/с$^2$.
Опорные значения плотности при 700~км и высоты масштаба 97~км
взяты из табличного ориентира NASA для периода около максимума
солнечной активности~\cite{atmosphere}. \emph{Одна} экспонента
на 500--1200~км является исследовательским приближением без
пространственно-временной изменчивости, F10.7 и Ap.

Высшие гармоники поля Земли, притяжение Солнца и Луны, световое
давление, тяга и коррекции орбит в уравнение~\eqref{eq:motion}
\emph{не включены}. Их возможный вклад на многих витках здесь
не оценён. Для согласованного расширения потребуются коэффициенты
поля, эфемериды, система времени и параметры поверхностей КА.

Поиск задаёт на общей эпохе отдельное состояние
$(\mathbf r_j(0),\mathbf v_j(0))$ каждого КА. Положения выбираются
по лучам целевой проекции с индивидуальной дальностью; скорости
задаются производной движущихся лучей и параметризованной
компонентой вдоль луча. Небольшое смещение внутренних начальных
узлов допускается, но крайние четыре угла остаются точными.
Ограниченный обзор стартового случая перебирает 12 комбинаций
трёх параметров. Далее состояние каждого аппарата преобразуется
в собственные классические \emph{оскулирующие} элементы
$(a,e,i,\Omega,\omega,\nu)$:
\begin{equation}
 \mathbf h=\mathbf r\times\mathbf v,\qquad
 \mathbf e=\frac{\mathbf v\times\mathbf h}{\mu}-\frac{\mathbf r}{r},
 \qquad e=|\mathbf e|,\qquad
 a=\left(\frac2r-\frac{v^2}{\mu}\right)^{-1}.
 \label{eq:elements}
\end{equation}
Углы $i,\Omega,\omega,\nu$ получаются из направлений
$\mathbf h$, линии узлов, вектора эксцентриситета и $\mathbf r$;
для кругового случая принято $\omega=0$, $\nu$ означает аргумент
широты. При обратном преобразовании используется
$p=a(1-e^2)$, $\mathbf r_{\rm pqw}=
\frac{p}{1+e\cos\nu}(\cos\nu,\sin\nu,0)^T$ и
$\mathbf v_{\rm pqw}=\sqrt{\mu/p}(-\sin\nu,e+\cos\nu,0)^T$,
затем поворот $R_3(\Omega)R_1(i)R_3(\omega)$. Совпадение исходных
и восстановленных состояний проверяется численно. Полный перечень
100 элементов с эпохой и конвенцией хранится в
\url{../orbit_study/drag_demo/results/case.json}; включать
его в основную таблицу отчёта нецелесообразно.

\section{Наблюдаемая геометрия и критерии}

Пусть $\mathbf s=(\mathbf r_0-\mathbf R)/D$,
$D=|\mathbf r_0-\mathbf R|$ --- единичное направление на КА~0.
Ортонормированные оси $\mathbf e_\xi,\mathbf e_\eta$ картинной
плоскости строятся из поперечной локальной оси $\mathbf e_y$ и
$\mathbf s$: первая есть нормированная проекция $\mathbf e_y$
на плоскость, перпендикулярную $\mathbf s$, вторая равна
$\mathbf s\times\mathbf e_\xi$. Для $\mathbf d_j=\mathbf r_j-\mathbf R$
гномоническая проекция определена при
$\mathbf d_j\cdot\mathbf s>10^{-9}$~км и равна
\begin{equation}
 \xi_j=\frac{\mathbf d_j\cdot\mathbf e_\xi}
             {\mathbf d_j\cdot\mathbf s},\qquad
 \eta_j=\frac{\mathbf d_j\cdot\mathbf e_\eta}
             {\mathbf d_j\cdot\mathbf s}.
 \label{eq:gnomonic}
\end{equation}
Координаты $(\xi,\eta)$ безразмерны. Подписи осей окна
$\arctan\xi$ и $\arctan\eta$ даны в градусах, но расстояние
между точками на этой плоскости не является истинным углом между
соответствующими лучами.

\emph{Заданная внешняя угловая сторона} $\theta=20^\circ$
означает длину $\alpha=\tan\theta$ каждой стороны \emph{целевого
внешнего контура} на гномонической плоскости. Целевой узел строки
$r$ и столбца $c$ имеет
$(\xi^*_{rc},\eta^*_{rc})=
(\alpha c/(N-1),\alpha r/(N-1))$.
Это не угловой шаг соседей. Если $b_k$ --- четыре фактические
длины сторон внешнего контура на картинной плоскости, то
\begin{equation}
 \theta_{\rm ext}=\arctan\!\left(\frac14\sum_{k=0}^3b_k\right).
 \label{eq:outer}
\end{equation}
Размер $\theta_{\rm ext}$ показывается отдельно от ошибок формы;
сами допуски к форме размер не закрепляют.

Для каждого КА проверяется положительная определённость
\eqref{eq:gnomonic} и угол места
\begin{equation}
 E_j=\arcsin\!\left(
 \frac{\mathbf d_j}{|\mathbf d_j|}\cdot\frac{\mathbf R}{R_E}
 \right)\geq15^\circ.
 \label{eq:elevation}
\end{equation}
Внешний контур должен быть выпуклым и иметь положительный порядок
вершин. При длинах его сторон $b_k$, диагоналях $d_1,d_2$ и
внутренних углах $\gamma_k$ применены допуски
\begin{equation}
\begin{gathered}
 \varepsilon_b=\max_k\frac{|b_k-\bar b|}{\bar b}\leq0{,}1,
 \qquad \bar b=\frac14\sum_kb_k,\\
 \varepsilon_d=\frac{2|d_1-d_2|}{d_1+d_2}\leq0{,}1,
 \qquad
 \varepsilon_\gamma=\max_k|\gamma_k-\pi/2|\leq10^\circ.
\end{gathered}
 \label{eq:shape}
\end{equation}
При $N>2$ дополнительно проверяются порядок всех строк и столбцов,
выпуклость каждой ячейки, ошибки её сторон, диагоналей и углов с
теми же пределами. Наилучшее подобие регулярной решётке
$z_j=\xi_j+\mathrm i\eta_j\approx\bar z+\kappa(q_j-\bar q)$,
$q_{rc}=c/(N-1)+\mathrm i r/(N-1)$, определяет ошибку узла
\begin{equation}
 \varepsilon_{\rm node}=
 \frac{\max_j|z_j-\bar z-\kappa(q_j-\bar q)|}{|\kappa|/(N-1)}
 \leq0{,}1,
 \quad
 \varepsilon_{\rm step}=
 \max_{(j,k)\in\mathcal E}
 \frac{|\,|z_j-z_k|-\bar l\,|}{\bar l}\leq0{,}1,
 \label{eq:grid}
\end{equation}
где $\mathcal E$ --- горизонтальные и вертикальные соседние пары,
$\bar l$ --- средняя длина их шага на плоскости. Для каждого такого
ребра проверяется также \emph{истинный} угол между направлениями
\begin{equation}
 \delta_{jk}=\arccos\!\left(
 \frac{\mathbf d_j\cdot\mathbf d_k}
 {|\mathbf d_j|\,|\mathbf d_k|}\right).
 \label{eq:neighbor}
\end{equation}
Порог $\min_{\mathcal E}\delta_{jk}$ выбирается отдельно как
$1^\circ$ либо $2^\circ$. Видимый квадрат принимается только там,
где выполнены условия для \emph{всех} $S$ аппаратов одновременно.
Связный интервал находится по соседним принятым временным узлам;
его длительность не включает неизвестное поведение за концами узлов.

Физические размеры вычисляются только из абсолютных радиусов:
$d_{jk}(t)=|\mathbf r_j(t)-\mathbf r_k(t)|$ в километрах.
Программа отдельно сообщает минимум и максимум соседних
пространственных сторон, четыре стороны и две диагонали между
крайними КА, максимум любой пары и минимум \emph{всех}
$S(S-1)/2$ пар. Условный порог сближения
$d_{\rm crit}=0{,}02+1{,}00=1{,}02$~км состоит из принятой суммы
радиусов КА и дополнительного запаса. Реальные габариты и
неопределённость состояний не дают оснований объявлять конфигурацию
безусловно безопасной.

\section{Стартовый случай: полученное и непроверенное}

Для $N=10$, $\theta=20^\circ$, порога истинного соседнего угла
$2^\circ$ и окна $t=-450\ldots450$~с найден принимаемый интервал
$-15\ldots15$~с (30~с) при шаге экрана 5~с.
Проверка проекции с шагом 1~с дала $-16\ldots17$~с (33~с).
На последнем интервале $\theta_{\rm ext}=19{,}6381\ldots
20{,}0000^\circ$, $\min\delta_{jk}=2{,}00156^\circ$,
$\min E_j=54{,}778^\circ$. Максимальные
$(\varepsilon_b,\varepsilon_d,\varepsilon_\gamma)$ равны
$(0{,}005331,0{,}001100,0{,}003951\,{\rm рад})$,
$\varepsilon_{\rm node}=0{,}08979$ шага.

На всём показанном окне фактические высоты КА составляют
568{,}36--1170{,}97~км. Все 4950 пар проверены на сетке 1~с;
межузловые минимумы найдены из смены знака
$(\mathbf r_j-\mathbf r_k)\cdot(\mathbf v_j-\mathbf v_k)$
с кубическим интерполянтом Эрмита. Ближайшее событие и кандидаты
ниже условного порога уточняются повторным интегрированием с шагом
не более 0{,}05~с. Минимум 2{,}458856~км относится к паре
9--10 при $t=-334{,}924205$~с: 5{,}8813$^\circ$~с.~ш. и
55{,}8144$^\circ$~в.~д. Это место получено из той же
приближённой ориентации Земли, что и проекция. В пределах окна
порог 1{,}02~км не нарушен.

При отключении сопротивления для \emph{тех же} начальных
элементов на окне максимальное расхождение положений составило
$4{,}05\cdot10^{-5}$~км; длительность при шаге 1~с осталась
33~с. Это лишь короткая проверка чувствительности. Оценки
оскулирующих перигеев и апогеев вне окна выводят часть аппаратов
примерно к 4740~км, за пределы принятой области плотности.
Ни полный виток, ни 30 суток с новой моделью не валидированы;
повторение квадрата и безопасность вне окна неизвестны.

Кнопка отдельного долгосрочного отчёта записывает выбранный горизонт,
элементы всех КА и причину, по которой \emph{долгосрочные численные
результаты не вычислены}. Отчёт создаётся только по нажатию;
графики последующих пролётов в основном окне отсутствуют.
Прежние результаты только с $J_2$ хранятся для истории проекта
и не подставляются как результат новой модели.

Пункт интерфейса ``С управлением'' имеется, но совместимого
рассчитанного решения \emph{нет ни для одного} $N=2,\ldots,10$.
Исторические данные относятся лишь к четырём КА ($N=2$), имеют
приложенное ускорение и реконструируемый видимый интервал,
но другую внешнюю сторону 1{,}8183$^\circ$, не имеют согласованной
UTC-эпохи и подтверждённого двигателя рассматриваемой платформы.

\section*{Источники данных и файлы воспроизведения}
\begin{thebibliography}{9}
\bibitem{atmosphere} NASA NTRS, табличная атмосферная плотность
  около максимума солнечной активности,
  \url{https://ntrs.nasa.gov/api/citations/19650024643/downloads/19650024643.pdf}.
\bibitem{corvus} Astro Digital, \emph{Landmapper Constellation Orbital Debris
  Assessment Report}, описание Corvus-HD 16U,
  \url{https://fcc.report/IBFS/SAT-LOA-20170508-00071/1228940.pdf}.
\end{thebibliography}

Настройки расчёта: \url{../orbit_study/drag_demo/config.json};
элементы и результаты одного пролёта:
\url{../orbit_study/drag_demo/results/case.json};
минимумы всех пар:
\url{../orbit_study/drag_demo/results/pair_minima.csv}.
Математический исходный черновик:
\url{inverse_projection_4sat_draft.tex}.

\end{document}
